\section{Introduction}\label{sec:intro}

Systems that answer questions over PDFs first extract the text, then cut it
into chunks that are indexed and retrieved. Because documents differ in layout
and structure, a natural idea is to choose the parser and chunker for each
document, and recent work proposes doing so from intrinsic quality
signals~\cite{demourajunior2026adaptive}. Before building such a chooser, we ask
how much retrieval quality is lost by using one fixed parse-and-chunk setup for
every PDF, and whether a choice per document, made by a cross-fitted oracle from
a part of each document's own questions, recovers any of it. The oracle selects
on few questions, so its gap is not a bound on every chooser. We study five
questions. \textbf{RQ1:} Does choosing a setup per document improve retrieval over
the best single setup? \textbf{RQ2:} At a fixed token budget, how does chunk size
relate to recall? \textbf{RQ3:} Does the best chunk size change when a fixed number
of chunks replaces the fixed token budget? \textbf{RQ4:} Do differences in answer
quality follow differences in retrieval? \textbf{RQ5:} Do tuned parser settings and
chunk sizes beat their defaults (a library default for the parsers and
LangChain, our own starting size for the other chunkers)?

Every setup reads its ranked chunks up to the same token budget. The RQ1 design,
the scorer and the RQ1 decision rule were written into dated specifications
before the first held-out set was scored; the designs of the other parts came
later and are dated in Section~\ref{sec:method}. We find no evidence of a gain from
choosing a setup per document, and the NO-GO decision is repeated on a second
held-out set with a different source of gold evidence. 
